\subsection{\rev{Positioning of this Review}}
\label{sec:review_positioning}
\begin{revision}
The closest comparisons are reviews of clinical gait analysis and reviews of the agreement between markerless and marker-based systems. Sethi et al. \cite{sethi2022comprehensive} review gait parameters and measurement methods, including the role of pose estimation in clinical gait analysis. Scataglini et al. \cite{scataglini2024accuracy} compare the accuracy, validity and reliability of three-dimensional gait measurement systems. Hulleck et al. \cite{hulleck2022present} discuss how clinical gait measures and scales are used, while Debnath et al. \cite{debnath2022review} examine computer vision across a wider range of rehabilitation and movement assessment tasks.

This review brings those questions together at the level of a joint and an intended use. It first asks which gait variables matter for hip, knee and ankle mobility, including changes at neighbouring joints. It then asks whether a video method measures the variable, recognises a pattern or supports an inference about a patient's condition. For example, a study that checks knee flexion angles against motion capture can inform measurement in the settings tested. Estimating the severity of knee impairment also requires evidence from patients assessed with a relevant clinical measure. Table~\ref{ExistingSurvey} shows where this emphasis overlaps with earlier reviews and what it adds for joint mobility assessment.
\end{revision}

\begin{table}[!ht]
\centering\small
\caption{\rev{Positioning relative to representative reviews of gait, measurement and movement analysis.}}
\label{ExistingSurvey}
\color{RevisionBlue}\hypersetup{allcolors=RevisionBlue}
\renewcommand{\arraystretch}{1.07}
\begin{tabularx}{0.98\linewidth}{@{}>{\raggedright\arraybackslash}p{0.24\linewidth}>{\raggedright\arraybackslash}p{0.31\linewidth}>{\raggedright\arraybackslash}X@{}}
\toprule
\textbf{Review} & \textbf{Main question or scope} & \textbf{How the present review uses or extends it} \\
\midrule
Sethi et al. \cite{sethi2022comprehensive} & Gait parameters, marker-based and markerless analysis, and pose estimation in clinical gait. & Follows a particular hip, knee or ankle variable from a clinical observation to its video estimate and interpretation. \\ \addlinespace[3pt]
Scataglini et al. \cite{scataglini2024accuracy} & Accuracy, validity and repeatability of three-dimensional markerless gait systems against marker-based systems. & Uses measurement evidence to judge individual variables and extends the comparison to video-based pattern recognition and clinical assessment. \\ \addlinespace[3pt]
Prakash et al. \cite{prakash2018recent} and Hulleck et al. \cite{hulleck2022present} & Gait assessment technologies and applications, including clinical measures, pathology and scales. & Uses the clinical findings to identify which variables and reference measurements matter for a particular joint and task. \\ \addlinespace[3pt]
Debnath et al. \cite{debnath2022review} & Computer vision for physical rehabilitation and movement assessment. & Examines how visual measures of walking relate to lower-limb joint impairment and compensation. \\ \addlinespace[3pt]
Lee et al. \cite{lee2014comprehensive}, Singh et al. \cite{singh2018vision}, Connor and Ross \cite{connor2018biometric}, Isaac et al. \cite{isaac2019trait} (preprint), Khan et al. \cite{khan2021vision}, Sepas-Moghaddam and Etemad \cite{sepas2022deep} & Gait sensing and representations for identifying people. & Draws on representation methods and evaluates their relevance to an explicitly defined mobility variable or outcome. \\ \addlinespace[3pt]
Saboor et al. \cite{saboor2020latest} & Wearable sensing and machine learning for gait analysis. & Places wearable measurements alongside video where they supply a reference or additional information about movement. \\ \addlinespace[3pt]
Kong and Fu \cite{kong2022human} and Gaya-Morey et al. \cite{gaya2024deep} & Action recognition and prediction, activity monitoring and fall detection. & Examines which methods can transfer to joint assessment and what validation of a gait pattern or clinical outcome requires. \\
\bottomrule
\end{tabularx}
\end{table}

\begin{revision}
The contribution is a joint-specific reading of the evidence. The review connects clinical gait findings to candidate video measures in Table~\ref{tab:joint_targets} and shows how the validation requirements change with the intended claim in Table~\ref{tab:validation_claims}. This helps readers judge when a study supports measurement, when it supports recognition of a pattern, and when a clinical conclusion still needs to be tested.
\end{revision}
\FloatBarrier
